% =====================================================================
%  TCET CODING PLATFORM -- Official Technical Documentation
%  MODULE 4 of 7 : Problems, Submissions & Code Execution Engine
%  STANDALONE FILE.  Compile: pdflatex Module4_Execution.tex  (twice)
% =====================================================================
\documentclass[12pt,a4paper,oneside]{report}
\usepackage[T1]{fontenc}
\usepackage[utf8]{inputenc}
\usepackage[a4paper,margin=1in]{geometry}
\usepackage{mathptmx}\usepackage{microtype}\usepackage{graphicx}\usepackage{xcolor}
\usepackage{tabularx}\usepackage{booktabs}\usepackage{longtable}\usepackage{xltabular}
\usepackage{array}\usepackage{enumitem}\usepackage{listings}\usepackage{titlesec}
\usepackage{fancyhdr}\usepackage{parskip}\usepackage[hidelinks]{hyperref}
\definecolor{tcetblue}{RGB}{0,0,0}\definecolor{tcetaccent}{RGB}{0,0,0}
\definecolor{codebg}{RGB}{245,246,248}\definecolor{codeframe}{RGB}{215,219,225}
\definecolor{codekw}{RGB}{0,0,0}\definecolor{codecomment}{RGB}{0,0,0}
\definecolor{codestr}{RGB}{0,0,0}\definecolor{boxframe}{RGB}{0,0,0}
\hypersetup{colorlinks=true,linkcolor=tcetblue,urlcolor=tcetaccent,citecolor=tcetblue,
  pdftitle={TCET Coding Platform -- Module 4: Execution Engine},pdfauthor={Centre of Excellence, TCET}}
\lstdefinestyle{tcp}{backgroundcolor=\color{codebg},basicstyle=\ttfamily\footnotesize,
  keywordstyle=\color{codekw}\bfseries,commentstyle=\color{codecomment}\itshape,
  stringstyle=\color{codestr},numberstyle=\tiny\color{gray},frame=single,
  rulecolor=\color{codeframe},breaklines=true,showstringspaces=false,tabsize=2,
  columns=fullflexible,keepspaces=true,xleftmargin=0.6em,framexleftmargin=0.6em}
\lstset{style=tcp}
\newcommand{\code}[1]{\texttt{#1}}
\newenvironment{callout}[1]{\par\vspace{4pt}\noindent
  \begin{tabular}{@{}p{0.014\linewidth}@{\hspace{0.6em}}p{0.93\linewidth}@{}}
  \textcolor{boxframe}{\rule{2pt}{\baselineskip}} & \textbf{\textcolor{tcetblue}{#1}}\\[2pt] & }{\end{tabular}\par\vspace{6pt}}
\titleformat{\chapter}[display]{\bfseries\fontsize{18}{22}\selectfont}{\bfseries\fontsize{14}{17}\selectfont Chapter \thechapter}{0.2em}{}
\titlespacing*{\chapter}{0pt}{6pt}{22pt}
\titleformat{\section}{\bfseries\fontsize{16}{19}\selectfont}{\thesection}{0.6em}{}
\titleformat{\subsection}{\bfseries\fontsize{14}{17}\selectfont}{\thesubsection}{0.5em}{}
\pagestyle{fancy}\fancyhf{}\renewcommand{\headrulewidth}{0.4pt}
\fancyhead[L]{\small\itshape TCET Coding Platform}\fancyhead[R]{\small\itshape Module 4 \textbullet\ Execution}
\fancyfoot[C]{\thepage}
\fancypagestyle{plain}{\fancyhf{}\fancyfoot[C]{\thepage}\renewcommand{\headrulewidth}{0pt}}
\newcommand{\note}[1]{\par\vspace{3pt}\noindent{\color{tcetaccent}\textbf{Note.}}~#1\par\vspace{3pt}}
\begin{document}
\begin{titlepage}\centering
\vspace*{1.4cm}{\color{tcetaccent}\rule{\linewidth}{1.2pt}}\\[0.5cm]
{\LARGE\bfseries TCET Coding Platform\par}
\vspace{0.4cm}{\Large Official Technical Documentation\par}{\large Draft~1 \textemdash\ Testing Phase\par}
{\color{tcetaccent}\rule{\linewidth}{1.2pt}}\\[1.4cm]
{\Huge\bfseries\color{tcetblue} Module 4\par}\vspace{0.3cm}
{\Huge Problems, Submissions \&\\[0.2cm] Code Execution Engine\par}\vspace{2.0cm}
{\large\bfseries Powered by\par}{\Large\color{tcetblue} Centre of Excellence (CoE), TCET\par}\vfill
\begin{tabular}{rl}\textbf{Module} & 4 of 7\\[3pt]
\textbf{Scope} & Problem bank, judging pipeline, sandbox, wrapping\\[3pt]
\textbf{Status} & Draft~1 (living document until Stage~2)\\[3pt]
\textbf{Audience} & Engineering, Officials\\\end{tabular}
\vspace{0.8cm}{\small\itshape \today\par}\end{titlepage}
\pagenumbering{roman}\tableofcontents\clearpage\pagenumbering{arabic}

% =====================================================================
\chapter{Overview}
% =====================================================================
\section{Scope of this Module}
This module specifies the core practice loop: how problems are authored and made
visible, how a submission travels from the editor to a graded verdict, and how
the sandboxed execution engine compiles and judges untrusted code across many
languages. It also specifies how ratings and statistics are derived and how the
pipeline recovers from failure.

\section{Key Concepts}
\begin{longtable}{@{}p{0.24\linewidth} p{0.68\linewidth}@{}}
\toprule
\textbf{Concept} & \textbf{Meaning}\\
\midrule
\endhead
Problem & An authored programming task with statement, limits, and tests.\\
Sample test & A visible test used by Run for immediate feedback.\\
Hidden test & A non-visible test used together with samples to judge Submit.\\
Run & Synchronous execution against samples only; unscored.\\
Submit & Asynchronous, queued execution against all tests; graded.\\
Verdict & The normalized outcome status of an execution.\\
Finalization & The one-time application of a verdict and its side effects.\\
\bottomrule
\end{longtable}

% =====================================================================
\chapter{The Problem Bank}
% =====================================================================
\section{Problem Lifecycle}
Every problem moves through a three-state lifecycle controlled by faculty.
\begin{longtable}{@{}p{0.20\linewidth} p{0.72\linewidth}@{}}
\toprule
\textbf{State} & \textbf{Meaning}\\
\midrule
\endhead
Draft & Authored but invisible to students; freely editable.\\
Published & Live and visible to eligible students.\\
Archived & Retired; retained for history but not offered to students.\\
\bottomrule
\end{longtable}

\section{Visibility Rules}
A student sees a problem only if it is Published and either open to all
departments or targeted at the student's own department. Faculty management views
expose problems in any state to their owners. These rules are enforced in the
service layer, not merely in the UI.

\section{Problem Definition}
A problem carries a rich statement (statement, input format, output format,
constraints), a difficulty (Easy/Medium/Hard), tags, and per-problem resource
limits (CPU time and memory). It defines two sets of test cases: \emph{sample}
tests shown to students and used by Run, and \emph{hidden} tests used together
with samples to judge a Submit.

\begin{callout}{Why sample and hidden tests are separated}
Sample tests give students fast, honest feedback while they work. Hidden tests
prevent a solution from being hard-coded to the visible cases, so a passing
Submit reflects a genuinely general solution.
\end{callout}

% =====================================================================
\chapter{Submission Modes}
% =====================================================================
\section{Run vs. Submit}
\begin{longtable}{@{}p{0.14\linewidth} p{0.78\linewidth}@{}}
\toprule
\textbf{Mode} & \textbf{Behaviour}\\
\midrule
\endhead
Run & Synchronous. Executes against sample tests and returns immediate feedback. Rate limited. Awards nothing.\\
Submit & Asynchronous. Persists a submission, enqueues it, and returns a receipt. A worker later runs the code against all tests and finalizes the verdict.\\
\bottomrule
\end{longtable}

\section{The Submit Pipeline}
\begin{lstlisting}[language={}]
1. Client sends code + language + problemId.
2. Service validates problem visibility; ensures the user record exists.
3. A submission is persisted with status QUEUED and the full test count.
4. The submission id is enqueued; the job id is stored.
   (If enqueue fails, the submission is marked errored immediately.)
5. Client receives a receipt and polls for the result.
6. A worker picks up the job -> status RUNNING.
7. The execution provider runs the code against all tests.
8. Finalization writes the verdict and, for practice submissions, recomputes
   user and problem aggregates and the leaderboard.
\end{lstlisting}

\section{Idempotent Finalization}
Finalization is guarded so a verdict is applied at most once: a submission that
already carries a finalization marker and a terminal status is returned
unchanged. This makes the pipeline safe against duplicate delivery, retries, and
the stale-recovery re-enqueue path.

\section{Polling Model}
Because Submit is asynchronous, the client polls the submission endpoint until a
terminal status is observed. The server returns the current status on each poll;
the verdict, runtime, memory, pass counts, and diagnostics become available once
judging finalizes.

% =====================================================================
\chapter{Rating \& Statistics}
% =====================================================================
\section{First-Solve Rating}
Rating points are awarded only for the \emph{first} accepted submission per
problem, weighted by difficulty. Re-solving an already-solved problem awards
nothing further. This rewards breadth of solved problems rather than repeated
submissions.

\section{Recomputed Aggregates}
On finalizing a practice submission the platform recomputes, from the user's
submissions and the problem's submissions:
\begin{itemize}[leftmargin=1.5em]
  \item \textbf{User:} rating, score, problems solved, submission count, accepted
        count, accuracy, and last-accepted time.
  \item \textbf{Problem:} total submissions, accepted submissions, and acceptance
        rate.
  \item \textbf{Leaderboard:} a denormalized entry for fast ranking.
\end{itemize}

\note{Only practice submissions contribute to a user's public rating and
accuracy. Contest coding submissions are graded within the contest and do not
inflate the global practice rating.}

\section{Analytics Derived}
From the same submission history the platform derives per-user analytics:
difficulty breakdown of solved problems, language-usage breakdown, an activity
heatmap over time, and recent accepted activity. These power the student profile
and the faculty view of a student.

% =====================================================================
\chapter{The Execution Engine}
% =====================================================================
\section{Provider Abstraction}
Execution is abstracted behind an execution-provider interface with two
operations: run against samples and execute against all tests. Two
implementations exist.
\begin{longtable}{@{}p{0.26\linewidth} p{0.66\linewidth}@{}}
\toprule
\textbf{Provider} & \textbf{Use}\\
\midrule
\endhead
Judge0 provider & Production; delegates to a Judge0 instance.\\
Stub provider & Deterministic stand-in for tests and offline development.\\
\bottomrule
\end{longtable}

\section{Sandboxing}
Judge0 runs each execution inside an \code{isolate} sandbox with \emph{no network
access} and strict CPU-time and memory limits. Untrusted student code therefore
cannot reach the network, exceed its resource budget, or affect the host.

\begin{callout}{Containment first}
The execution engine treats every submission as hostile. Network is disabled,
resources are bounded, and each run is isolated, so even deliberately malicious
code is confined to a disposable sandbox.
\end{callout}

\section{Language Resolution}
The provider resolves the concrete language runtime dynamically: it queries the
live language list and selects a preferred version by name, then by name hints,
then by a curated fallback identifier. Separate fallback sets exist for a local
and a cloud execution backend, so the correct runtime is chosen in either
deployment.

\section{Fair Limits Across Languages}
Interpreted and virtual-machine languages are inherently slower and heavier than
compiled ones. The provider applies per-language CPU-time and memory multipliers
so that, for example, an interpreted language is given proportionally more
wall-clock and memory than C or C++, subject to an absolute ceiling. This yields
fair judging across a diverse language set.

\section{Chunked Judging with Short-Circuit}
For a full submission the provider runs a compilation preflight and then executes
test cases in small parallel chunks, stopping as soon as a non-accepted result
appears. A failing solution is therefore not run against every remaining test,
which conserves execution capacity under load.

\section{Verdict Normalization \& Aggregation}
Raw execution outcomes are mapped to a normalized status set, and a multi-test
submission is reduced to a single verdict by priority.
\begin{longtable}{@{}p{0.30\linewidth} p{0.62\linewidth}@{}}
\toprule
\textbf{Status} & \textbf{Meaning}\\
\midrule
\endhead
Accepted & All executed tests passed.\\
Wrong answer & Output did not match on at least one test.\\
Time-limit exceeded & A test exceeded its time budget.\\
Compilation error & The program failed to compile.\\
Runtime error & The program crashed or exited abnormally.\\
Internal error & An infrastructure failure occurred.\\
\bottomrule
\end{longtable}

\noindent The final status is chosen by priority (compilation error dominates,
then runtime error, then time-limit, then wrong answer, then accepted), with
internal error dominating if any infrastructure failure occurred; the most
informative diagnostic is surfaced to the user, along with aggregate pass count,
worst-case runtime, and peak memory.

% =====================================================================
\chapter{Language Support \& Automatic Wrapping}
% =====================================================================
\section{Supported Languages}
The platform supports 20+ executable languages (including C, C++, Java, Python,
JavaScript, TypeScript, Go, Rust, C\#, PHP, Ruby, Kotlin, Swift, Dart, Scala,
Elixir, Erlang, and Assembly) plus three editor-only languages (HTML, CSS, React)
that are never sent to the judge. Aliases exist \textemdash\ for instance Arduino
compiles as C++, and a ``vanilla'' variant runs as JavaScript.

\section{The Code Wrapper}
A distinctive feature lets students write function-style, class-based solutions
without a main entry point. When a submission defines a solution class but no
entry point, the platform synthesizes one:
\begin{itemize}[leftmargin=1.5em]
  \item It detects the target method by scanning the class body, ignoring
        constructors and reserved names.
  \item It generates an entry point that instantiates the solution, invokes the
        method, and prints the result using a language-specific formatter that
        renders booleans, arrays/lists, pairs, and maps consistently.
  \item Wrappers exist for the major languages; others pass through unchanged.
\end{itemize}
If the student already provides an entry point, their program is used verbatim.
This dual behaviour supports both function-style and full-program problems from
the same editor, and it behaves identically for practice and contest submissions.

\begin{callout}{Consistent output formatting}
Because the generated wrapper formats results uniformly across languages, the
same logical answer prints the same way whether the student wrote Python, Java,
or C++, so a single expected-output test judges all languages fairly.
\end{callout}

\section{Editor-Only Languages}
HTML, CSS, and React exist for authoring and preview contexts and are explicitly
blocked from execution: the engine refuses to send them to the judge, so they can
never consume sandbox capacity or be mistaken for executable submissions.

% =====================================================================
\chapter{Failure Handling \& Recovery}
% =====================================================================
\section{Failure Modes \& Responses}
\begin{longtable}{@{}p{0.30\linewidth} p{0.62\linewidth}@{}}
\toprule
\textbf{Failure} & \textbf{Response}\\
\midrule
\endhead
Enqueue failure & The submission is finalized as an internal error immediately, never left pending.\\
Worker exception & The submission is finalized as an internal error rather than left running.\\
Judge transient error & Retried automatically with backoff via the queue.\\
Worker crash mid-job & Stale recovery re-enqueues the stuck submission.\\
\bottomrule
\end{longtable}

\section{Stale-Submission Recovery}
Submissions left pending past a staleness threshold, and not yet finalized, are
re-enqueued by a recovery routine. This guards against a worker that died
mid-judging, so no student's submission is permanently stranded.

\section{Idempotency Guarantees}
Every recovery and retry path is safe because finalization is idempotent: a
submission that has already been finalized is returned unchanged rather than
re-judged, so aggregates and ratings cannot be double-counted.

% =====================================================================
\chapter{End-to-End Walkthroughs}
% =====================================================================
\section{Walkthrough: Accepted Practice Submission}
\begin{enumerate}[leftmargin=1.6em]
  \item A student writes a solution and presses Run; sample tests pass.
  \item The student presses Submit; a submission is persisted and queued.
  \item A worker judges the submission against all tests; every test passes.
  \item The verdict is finalized as accepted; because it is the first acceptance
        of this problem, difficulty-weighted rating is awarded.
  \item The student's analytics and the leaderboard update.
\end{enumerate}

\section{Walkthrough: Failing Submission}
\begin{enumerate}[leftmargin=1.6em]
  \item The student submits a solution that fails a hidden test.
  \item The worker judges in chunks and stops at the first non-accepted test.
  \item The verdict is finalized as wrong answer with a diagnostic; no rating is
        awarded; accuracy reflects the attempt.
\end{enumerate}

\section{Walkthrough: Infrastructure Hiccup}
\begin{enumerate}[leftmargin=1.6em]
  \item A submission is queued but the worker crashes mid-judging.
  \item The submission remains pending past the staleness threshold.
  \item Stale recovery re-enqueues it; a healthy worker judges it to completion.
  \item Finalization applies exactly once thanks to the idempotency guard.
\end{enumerate}

% =====================================================================
\chapter{Design Rationale \& Summary}
% =====================================================================
\section{Why this Design}
\begin{itemize}[leftmargin=1.5em]
  \item \textbf{Asynchial judging} keeps the platform responsive under mass load.
  \item \textbf{Sandboxing} makes running untrusted code safe by default.
  \item \textbf{Per-language limits} keep judging fair across runtimes.
  \item \textbf{Short-circuit judging} conserves capacity for failing solutions.
  \item \textbf{Idempotent finalization} makes retries and recovery safe.
  \item \textbf{The wrapper} lowers the barrier to entry for students while
        keeping judging consistent.
\end{itemize}

\section{Module Summary}
The execution engine turns a student's code into a fair, safe, and reproducible
verdict. Submissions are accepted instantly, judged asynchronously inside a
sandbox against sample and hidden tests, normalized to a consistent verdict,
and finalized exactly once, updating ratings and analytics for practice work.
The same engine underpins contest coding questions, which are governed by the
contest engine described in Module~5.

\vspace{0.6cm}
\begin{center}
\textit{--- End of Module 4. Continued in Module 5: Contest \& Assessment Engine. ---}
\end{center}
\end{document}
